\chapter*{Course Map}
\phantomsection%
\addcontentsline{toc}{chapter}{Course Map}

The course progresses from data organization to algorithm selection. Five stages build core skills; the final project compares algorithms with evidence.

\begin{mlfigure}
\begin{tikzpicture}[
  node distance=0.48cm and 0.72cm,
  box/.style={draw=ADSBlue,fill=ADSPaleBlue,rounded corners=2mm,text width=5.15cm,minimum height=0.92cm,align=left,inner xsep=4mm,font=\small\color{ADSNavy}},
  arr/.style={-{Stealth[length=2.2mm]},thick,draw=ADSTeal}
]
\node[box] (a) {\textbf{\color{ADSBlue}1. Foundations}\\[-0.15em]{\scriptsize Arrays and memory}};
\node[box,right=of a] (b) {\textbf{\color{ADSBlue}2. Linear structures}\\[-0.15em]{\scriptsize Lists, stacks, and queues}};
\node[box,below=of b] (c) {\textbf{\color{ADSBlue}3. Hierarchies}\\[-0.15em]{\scriptsize Trees, BSTs, and heaps}};
\node[box,left=of c] (d) {\textbf{\color{ADSBlue}4. Networks}\\[-0.15em]{\scriptsize Graph representation and traversal}};
\node[box,below=of d] (e) {\textbf{\color{ADSBlue}5. Core algorithms}\\[-0.15em]{\scriptsize Searching, sorting, and hashing}};
\node[box,right=of e] (f) {\textbf{\color{ADSBlue}6. Synthesis}\\[-0.15em]{\scriptsize Evidence-based comparison}};
\draw[arr] (a) -- (b);
\draw[arr] (b) -- (c);
\draw[arr] (c) -- (d);
\draw[arr] (d) -- (e);
\draw[arr] (e) -- (f);
\end{tikzpicture}
\end{mlfigure}

\begin{mltable}
\begin{tabularx}{\linewidth}{@{}>{\bfseries\leavevmode\color{ADSNavy}}p{0.13\linewidth}p{0.26\linewidth}X@{}}
\toprule
\rowcolor{ADSPaleBlue!92}
Stage & Main structures & Central question \\
\midrule
1 & Primitive/non-primitive structures, arrays & How is data represented and accessed in memory? \\
2 & Linked lists, circular lists, stacks, queues & How can pointer-based structures support dynamic or constrained access? \\
3 & Trees, binary trees, BSTs, heaps, balancing & How can hierarchy improve organization and searching? \\
4 & Graphs, adjacency structures, BFS/DFS, spanning trees & How can general relationships between objects be represented and traversed? \\
5 & Searching, sorting, hashing & How do algorithm choice and data organization affect running time? \\
6 & Final comparison project & How can two related algorithms be compared fairly using complexity, operation counts, runtime, and practical use cases? \\
\bottomrule
\end{tabularx}
\end{mltable}

\section*{How these notes use the course sources}
The notes preserve the terminology and topic order of the supplied slide decks, then use the official lab manual to strengthen objectives, outcomes, and problem sets. The final-project guideline is incorporated as a dedicated synthesis chapter. Slide/manual diagrams are redrawn as TikZ where useful, and source tasks are turned into C++ practice milestones with verification checklists.

\begin{clarification}
A few slide statements need context to avoid ambiguity in implementation. These notes do not silently replace them. Whenever an indexing convention, complexity claim, or standard definition needs extra context, the slide claim is stated first and the extra computer-science convention is placed in a box like this one.
\end{clarification}

\section*{Complexity vocabulary used across the course}
Big-O notation describes an asymptotic upper bound on how resource use grows with input size $n$; it is not an exact operation count. The table and chart use representative functions $f(n)$ to compare growth rates consistently.

\begin{mltable}
\begin{tabularx}{\linewidth}{@{}>{\bfseries\leavevmode\color{ADSNavy}}p{0.11\linewidth}p{0.16\linewidth}p{0.21\linewidth}X@{}}
\toprule
\rowcolor{ADSPaleBlue!92}
Class & Growth rate & Function plotted & Typical examples \\
\midrule
$O(1)$ & Constant & $f(n)=1$ & Array access; hash computation. \\
$O(\log n)$ & Logarithmic & $f(n)=\log_2 n$ & Binary search; balanced-BST lookup. \\
$O(n)$ & Linear & $f(n)=n$ & Linear search; linked-list traversal. \\
$O(n\log n)$ & Linearithmic & $f(n)=n\log_2 n$ & Merge sort; heap sort. \\
$O(n^2)$ & Quadratic & $f(n)=n^2$ & Bubble sort; selection sort. \\
$O(n!)$ & Factorial & $f(n)=n!$ & Permutation enumeration; brute-force traveling salesperson search. \\
\bottomrule
\end{tabularx}
\end{mltable}

\begin{mlfigure}
\centering
\begin{tikzpicture}
\begin{axis}[
  width=0.94\linewidth,
  height=7.6cm,
  domain=2:10,
  samples=120,
  xmin=2,xmax=12.7,
  ymin=0.8,ymax=10000000,
  ymode=log,
  log basis y=10,
  xtick={2,4,6,8,10},
  ytick={1,10,100,1000,10000,100000,1000000,10000000},
  xlabel={Input size, $n$},
  ylabel={Representative operations, $f(n)$ (log scale)},
  title={Representative growth by complexity class},
  title style={font=\bfseries\color{ADSNavy}},
  label style={font=\small\color{ADSNavy}},
  tick label style={font=\scriptsize\color{ADSMuted}},
  grid=major,
  grid style={draw=ADSPaleBlue},
  axis line style={draw=ADSMuted},
  axis x line*=bottom,
  axis y line*=left,
  clip=false
]
\addplot[very thick,draw=ADSMuted,densely dotted] {1};
\addplot[very thick,draw=ADSTeal,dashdotted] {ln(x)/ln(2)};
\addplot[very thick,draw=ADSBlue] {x};
\addplot[very thick,draw=ADSOrange,dashed] {x*ln(x)/ln(2)};
\addplot[very thick,draw=ADSRed,densely dashed] {x^2};
\addplot[very thick,draw=ADSNavy,mark=*,mark size=1.5pt] coordinates {
  (2,2) (3,6) (4,24) (5,120) (6,720) (7,5040) (8,40320) (9,362880) (10,3628800)
};
\node[anchor=west,font=\scriptsize\bfseries,text=ADSMuted] at (axis cs:10.15,1) {$O(1)$};
\node[anchor=west,font=\scriptsize\bfseries,text=ADSTeal] at (axis cs:10.15,3.322) {$O(\log n)$};
\node[anchor=west,font=\scriptsize\bfseries,text=ADSBlue] at (axis cs:10.15,10) {$O(n)$};
\node[anchor=west,font=\scriptsize\bfseries,text=ADSOrange] at (axis cs:10.15,33.22) {$O(n\log n)$};
\node[anchor=west,font=\scriptsize\bfseries,text=ADSRed] at (axis cs:10.15,100) {$O(n^2)$};
\node[anchor=west,font=\scriptsize\bfseries,text=ADSNavy] at (axis cs:10.15,3628800) {$O(n!)$};
\end{axis}
\end{tikzpicture}

\smallskip
{\small\color{ADSMuted}The logarithmic vertical scale keeps all six functions visible. Each curve follows the representative function listed in the table; actual running time depends on the implementation and hardware.}
\end{mlfigure}

\begin{keyidea}
Factorial complexity, $O(n!)$, arises when an algorithm examines every ordering of $n$ items. It becomes impractical very quickly: $10! = 3{,}628{,}800$ possible orderings, while $15!$ exceeds one trillion. Brute-force permutation problems and exhaustive solutions to the traveling salesperson problem are common examples.
\end{keyidea}
